\subsection{AdA can leverage prompting with first-person demonstrations}

To determine whether AdA can learn in zero-shot from first-person demonstrations, we prompted it with a first-person demonstration by a fine-tuned teacher, as follows. The teacher took control of the avatar in the first trial, while AdA continued to receive observations as usual, conditioning its Transformer memory. AdA was then allowed to proceed on its own for the remaining trials and its scores recorded in the usual manner. 

Figure \ref{fig:first-person-demonstration} shows the median score on our hand-authored test set of prompted AdA compared to an unprompted baseline. Prompted AdA is unable to exactly mimic the teacher's demonstration in the second trial of a median task, shown by a drop in score. It does, however, outperform an unprompted baseline across all numbers of trials, indicating that it is able to profitably incorporate information from the demonstration into its policy. This process is analogous to prompting of large language models, where the agent's memory is primed with an example of desired behaviour from which it continues. We note that AdA was never trained with such off-policy first-person demonstrations, yet its in-context learning algorithm is still able to generalise to these. 

In Figure \ref{fig:prompting-all-tasks} we provide prompting results for all single-agent hand-authored tasks and discuss the circumstances under which prompting is effective. In Appendix \ref{app:prompting} we also provide early results investigating prompting with human demonstrations on a subset of tasks. These reveal remarkable success in some cases, but also confirm that human demonstrations cannot  overcome inherent limitations of AdA's task distribution. Two videos compare the behaviour when \href{https://youtu.be/jdf4jkiJYa4}{prompted} and when \href{https://youtu.be/CuxPknjnYOI}{not prompted} on the task \texttt{Object permanence: yellow cube}.
